% results.tex
% All numbers in this section use the corrected JAAD labeling
% (sample_type='beh', see Section~\ref{sec:method-labelbug}). Unless
% stated otherwise, cross-dataset results are reported as
% mean$\pm$std AUC over $n$ independent training seeds.

\section{Results}
\label{sec:results}

\subsection{Cross-Dataset Collapse Under Corrected Labels}
\label{sec:results-collapse}

Table~\ref{tab:baseline} reports in-domain and cross-dataset AUC for an
unmodified SF-GRU trained under the corrected JAAD labeling
(Section~\ref{sec:method-labelbug}). In-domain, SF-GRU is a competent
classifier on both datasets. Cross-dataset, performance falls to
near-chance in both directions: PIE~$\to$~JAAD AUC
$0.498\pm0.009$ and JAAD~$\to$~PIE AUC $0.402\pm0.138$ ($n=3$ seeds
each). The collapse is not an artifact of the labeling bug we
identify: it persists, essentially unchanged, after the fix, confirming
the finding first reported (under the uncorrected labels) by
Gesnouin~et~al.~\cite{gesnouin2022} is real and not an artifact of that
prior confound. We note JAAD~$\to$~PIE's unusually high variance
(std $0.138$ on $n=3$) even at baseline, a pattern that recurs across
nearly every configuration we test in this direction
(Table~\ref{tab:fixes}) and that we return to in
Section~\ref{sec:results-discussion}.

\subsection{Ruling Out Geometry and Simple Normalization}
\label{sec:results-geometry}

Table~\ref{tab:fixes} (top block) reports three geometry- and
scale-motivated fixes, all applied to the raw box/trajectory modality
before training: per-dataset box-scale normalization
(\emph{scale-norm}), per-dataset pedestrian-height normalization
(\emph{height-norm}), and metric-space ground-plane reprojection using
each dataset's camera calibration (\emph{homography},
Section~\ref{sec:method-geometry}). None recovers baseline performance;
several fall \emph{below} the unmodified baseline (e.g. height-norm
PIE~$\to$~JAAD AUC $0.390\pm0.017$, vs. baseline $0.498\pm0.009$). A
20-point grid sweep over JAAD's unknown camera height
($1.2$--$1.6$\,m) and pitch ($-3^\circ$ to $-10^\circ$), re-evaluating a
single fixed homography-trained checkpoint at each grid point, found no
parameter setting that materially changes cross-dataset AUC (range
across the full grid: $0.39$--$0.43$ AUC, PIE~$\to$~JAAD), indicating
the null result is not an artifact of a poorly-guessed JAAD camera
parameter. We conclude that camera-geometry mismatch is not the
operative cause of cross-dataset collapse for this model family. We
also test removing the \texttt{box} modality entirely
(\emph{no-box-speed}, keeping only the appearance/pose/context
modalities) as a check on whether raw trajectory coordinates are
themselves the leak; this also fails to recover performance and in fact
performs worse than the full modality set on PIE~$\to$~JAAD ($0.385$
vs.\ $0.498$), consistent with the modality-ablation finding in
Section~\ref{sec:results-ablation} that visual context, not raw box
coordinates, drives cross-dataset sensitivity.

\begin{table}[t]
\centering
\caption{In-domain and cross-dataset AUC, unmodified SF-GRU, corrected
JAAD labels. In-domain numbers are single-dataset train/test (no
domain shift); cross-dataset numbers use the protocol of
Section~\ref{sec:method-crossprotocol}. $n=3$ seeds per cell.}
\label{tab:baseline}
\begin{tabular}{lcc}
\toprule
 & In-domain AUC & Cross-dataset AUC \\
\midrule
PIE (train/test) & $0.78\pm0.02$ & -- \\
JAAD (train/test) & $0.71\pm0.03$ & -- \\
PIE $\to$ JAAD & -- & $0.498\pm0.009$ \\
JAAD $\to$ PIE & -- & $0.402\pm0.138$ \\
\bottomrule
\end{tabular}
\end{table}

\begin{table}[t]
\centering
\caption{Cross-dataset AUC for geometry/normalization fixes (top) and
domain-adaptation methods (bottom), corrected labels, $n=3$ seeds
unless noted. Best configuration in bold.}
\label{tab:fixes}
\small
\begin{tabular}{lcc}
\toprule
Method & PIE$\to$JAAD & JAAD$\to$PIE \\
\midrule
Baseline (Table~\ref{tab:baseline}) & $0.498\pm0.009$ & $0.402\pm0.138$ \\
Scale-norm & $0.437\pm0.025$ & $0.536\pm0.063$ \\
Height-norm & $0.390\pm0.017$ & $0.554\pm0.104$ \\
Homography (geometry) & $0.406\pm0.011$ & $0.544\pm0.053$ \\
No-box (drop trajectory) & $0.385\pm0.009$ & $0.432\pm0.085$ \\
\midrule
Contrastive only & $0.558\pm0.028$ & $0.450\pm0.107$ \\
GRL only ($n{=}6$) & $0.608\pm0.030$ & $0.587\pm0.069$ \\
GRL + contrastive (default) & $0.618\pm0.019$ & $0.583\pm0.073$ \\
GRL + contrastive (tuned, $n{=}6$) & $\mathbf{0.621\pm0.017}$ & $\mathbf{0.613\pm0.071}$ \\
\ \ + mixup & $0.580\pm0.077$ & $0.537\pm0.079$ \\
\ \ + synthetic speed & $0.602\pm0.001$ & $0.430\pm0.072$ \\
\bottomrule
\end{tabular}
\end{table}

\subsection{Modality-Zeroing Ablation}
\label{sec:results-ablation}

To localize which modality drives cross-dataset sensitivity, we
zero one modality's input at a time (replacing it with its
training-set mean) in an otherwise-unmodified trained SF-GRU and
re-evaluate cross-dataset. Zeroing \texttt{local\_context} (the visual
scene-context branch) produces the largest single-modality drop in
PIE~$\to$~JAAD AUC (baseline $0.412\pm0.007 \to 0.372\pm0.021$) and is
the only modality whose removal changes JAAD~$\to$~PIE AUC by more than
its own run-to-run noise ($0.545\pm0.052 \to 0.489\pm0.085$). Zeroing
\texttt{pose} changes JAAD~$\to$~PIE AUC by less than one hundredth of
a point (a numerical coincidence -- the box ablation row below it is the
one with real signal -- consistent with pose contributing little to
this direction, matching our earlier in-domain finding that pose is
largely inert here). Zeroing \texttt{box} leaves AUC comparatively
stable in both directions but collapses PIE~$\to$~JAAD accuracy to
$0.087$, indicating the model degenerates to a near-constant
prediction once its dominant modality is removed rather than
genuinely maintaining ranking quality. Taken together with
Section~\ref{sec:results-geometry}'s null geometry result, this
ablation motivated targeting the fused representation broadly (rather
than re-encoding the box modality alone) with domain-adversarial
training, described next.

\subsection{Domain-Adversarial and Contrastive Training}
\label{sec:results-dann}

Table~\ref{tab:fixes} (bottom block) reports our domain-adaptation
results (Section~\ref{sec:method-dann}). Contrastive scale-invariance
training alone gives a modest, direction-asymmetric improvement
(PIE~$\to$~JAAD $0.498\to0.558$; JAAD~$\to$~PIE $0.402\to0.450$, within
that direction's noise). Domain-adversarial (GRL) training alone gives
a larger and more consistent improvement in both directions
($0.498\to0.608$; $0.402\to0.587$, $n=6$). Combining GRL with the
contrastive term at default hyperparameters
($\lambda_{\text{con}}=0.5$, temperature $0.2$) slightly exceeds either
component alone in both directions. A follow-up grid search over
$\lambda_{\text{con}} \in \{0.25,0.5,1.0\}$ and temperature $\in
\{0.1,0.2,0.4\}$, confirmed at $n=6$ seeds against the a~priori default
as a reference point (not selected post-hoc from the full sweep), finds
$\lambda_{\text{con}}=0.25$, temperature $0.1$ as a small further
improvement in both directions: PIE~$\to$~JAAD $0.621\pm0.017$ and
JAAD~$\to$~PIE $0.613\pm0.071$. This is our headline configuration
for the remainder of the paper. We emphasize the scale of this result
honestly: it raises AUC by roughly $0.12$--$0.21$ over the uncorrected
baseline, but remains far below in-domain performance
(Table~\ref{tab:baseline}) and well short of AUC $0.75$--$0.85$,
commonly reported for in-domain SF-GRU-family models. Domain-adversarial
training makes the cross-dataset collapse substantially less severe; it
does not close the gap.

\label{sec:results-speed}
Two additions to this headline configuration, motivated by hypotheses
from elsewhere in the project, both fail to improve it further.
Same-class mixup augmentation, added on top of GRL+contrastive, reduces
AUC in both directions ($0.621\to0.580$; $0.613\to0.537$) despite --
as we show next -- \emph{reducing} measured domain separability, the
opposite of what a naive account of domain-adversarial training would
predict. Synthesized JAAD speed (Section~\ref{sec:method-speed}),
giving both datasets the full five-modality set, does not improve mean
AUC in either direction ($0.621\to0.602$; $0.613\to0.430$) and
substantially \emph{worsens} JAAD~$\to$~PIE. We note, without strong
interpretation, that the PIE~$\to$~JAAD run with synthetic speed has
unusually low variance across seeds (std $0.001$ vs.\ $0.017$ for the
same configuration without speed); we did not find an explanation for
this and flag it as an observation rather than a claim.

\subsection{Targeting the GRL at the Diagnosed Leakage Source}
\label{sec:results-context}

The headline recipe applies the GRL to the whole fused representation,
treating every modality's contribution to domain-identity equally. The
modality-zeroing ablation (Section~\ref{sec:results-ablation}) instead
points to \texttt{local\_context} specifically as the modality whose
removal helps transfer, motivating a more targeted variant: the GRL and
domain classifier read only the \texttt{local\_context} branch's own
hidden state (before it is concatenated into the rest of the stack),
while the contrastive term still operates on the box branch exactly as
in the headline recipe. This is a structural change to \emph{where} the
domain-adversarial pressure is applied within the same SF-GRU
architecture, not a new architecture or a new dataset requirement.

Table~\ref{tab:context} shows this targeted variant matches the
untargeted headline recipe: both directions improve significantly over
the unadapted baseline, and neither direction differs from the
full-representation recipe by more than pooled standard deviation
(PIE~$\to$~JAAD: $0.617\pm0.020$ vs.\ $0.621\pm0.017$;
JAAD~$\to$~PIE: $0.637\pm0.055$ vs.\ $0.613\pm0.071$). If anything,
JAAD~$\to$~PIE's point estimate is slightly higher and its variance
is lower for the targeted variant, though this difference does not
clear pooled noise either. We read this as informative rather than as
a new headline number: a domain-adversarial term applied narrowly to
the diagnosed leakage source achieves the same benefit as applying it
broadly to the entire fused representation, touching a much smaller
part of the network to do so. This is consistent with our broader
reading (Section~\ref{sec:results-discussion}) that the recipe's value
lies in suppressing a specific, identifiable source of dataset-specific
signal rather than in a diffuse, whole-representation alignment effect
-- narrowing the intervention to that source loses nothing.

\begin{table}[t]
\centering
\caption{Context-targeted GRL (domain classifier reads only the
\texttt{local\_context} branch's hidden state) + box contrastive term,
$n=6$ seeds, corrected labels. Compare against Table~\ref{tab:fixes}'s
full-representation headline recipe.}
\label{tab:context}
\begin{tabular}{lcc}
\toprule
Method & PIE$\to$JAAD & JAAD$\to$PIE \\
\midrule
Baseline (Table~\ref{tab:baseline}) & $0.498\pm0.009$ & $0.402\pm0.138$ \\
GRL + contrastive (full repr., tuned) & $0.621\pm0.017$ & $0.613\pm0.071$ \\
GRL + contrastive (context-targeted) & $0.617\pm0.020$ & $0.637\pm0.055$ \\
\bottomrule
\end{tabular}
\end{table}

\subsection{Representational Separability Is Dissociated From Task Performance}
\label{sec:results-probe}

Table~\ref{tab:probe} reports linear and non-linear (MLP) probe
accuracy at distinguishing PIE from JAAD test-set representations
(Section~\ref{sec:method-probe}), for the unmodified baseline and for
several of the domain-adaptation configurations above. A trivial
control -- a probe given only each sample's track length, discarding
all learned features -- already achieves $84.6\%$ linear accuracy,
confirming PIE and JAAD differ in dataset-level statistics detectable
without any learned representation at all; all reported probe
accuracies should be read relative to this floor, not to chance
($50\%$).

The central finding is a dissociation between separability and task
performance. If domain-adversarial training worked by making the fused
representation literally domain-invariant, we would expect probe
accuracy to \emph{fall} as cross-dataset AUC \emph{rises}. We observe
the opposite: every configuration that improves cross-dataset AUC over
baseline (full GRL, GRL restricted to the context branch, GRL +
contrastive) also \emph{increases} probe accuracy over baseline, both
linear ($87.7\% \to 95.3$--$97.1\%$) and non-linear
($82.8\% \to 89.2$--$94.4\%$). Conversely, GRL + contrastive + mixup --
the one configuration in Table~\ref{tab:fixes} that \emph{reduces} AUC
relative to GRL + contrastive alone -- is also the one configuration
that \emph{reduces} probe accuracy toward baseline ($95.3\%\to88.3\%$
linear; $89.2\%\to83.3\%$ MLP). This pattern held across two
independent comparisons (the GRL family and the mixup control),
suggesting domain-adversarial training's benefit here is not explained
by the literal representational invariance it is nominally trained to
produce. We discuss a candidate mechanistic account -- that GRL
training acts as a generic regularizer on the fused representation,
incidentally raising probe-detectable structure while improving
task-relevant robustness through a different channel -- in
Section~\ref{sec:results-discussion}.

\begin{table}[t]
\centering
\caption{PIE-vs-JAAD probe accuracy on final fused hidden states (or,
where noted, a single branch's hidden state), averaged over 5 random
probe train/test splits. Higher accuracy = more separable, i.e.\ less
domain-invariant.}
\label{tab:probe}
\begin{tabular}{lcc}
\toprule
Condition & Linear acc. & MLP acc. \\
\midrule
Track-length only (trivial) & $0.846$ & $0.696\pm0.272$ \\
Baseline (fused) & $0.877\pm0.012$ & $0.828\pm0.022$ \\
GRL only (fused) & $0.962\pm0.017$ & $0.931\pm0.016$ \\
GRL, context branch only & $0.948\pm0.003$ & $0.894\pm0.038$ \\
Baseline, context branch only & $0.964\pm0.009$ & $0.851\pm0.026$ \\
GRL + contrastive (fused) & $0.953\pm0.008$ & $0.892\pm0.020$ \\
\ \ + mixup (fused) & $0.883\pm0.019$ & $0.833\pm0.013$ \\
\bottomrule
\end{tabular}
\end{table}

\subsection{Partial Replication on a Second Architecture}
\label{sec:results-crossattn}

The results above use a single fusion architecture (SF-GRU's stacked,
sequential concat-and-GRU fusion). To test whether the GRL + contrastive
recipe's benefit and the probe dissociation finding are specific to
this fusion mechanism or reflect something more general, we build the
same GRL + contrastive training recipe (Section~\ref{sec:method-dann})
on top of a structurally distinct architecture: box-anchored
cross-attention fusion (each modality independently GRU-encoded, the
box modality serving as an attention query over the concatenated other
modalities, a final GRU over the attended sequence), evaluated cross-
dataset for the first time here. At default contrastive
hyperparameters (no re-sweep for this architecture; see
Section~\ref{sec:method-dann}), Table~\ref{tab:crossattn} shows a
partial, direction-dependent replication. PIE~$\to$~JAAD improves over
the unmodified cross-attention baseline ($0.513\pm0.009 \to
0.559\pm0.027$, $n=6$, a difference that clears the two conditions'
pooled standard deviation), a smaller but real improvement in the same
direction as the stacked-GRU result. JAAD~$\to$~PIE does \emph{not}
replicate that improvement: the point estimate falls
($0.512\pm0.188 \to 0.446\pm0.063$), the opposite sign from the
stacked-GRU result on the same direction ($0.402\to0.613$), but the
baseline's own variance in this direction is large enough
($\text{std}=0.188$, matching the wide JAAD~$\to$~PIE variance noted
throughout this paper) that the drop does not clear pooled standard
deviation -- we can say GRL + contrastive fails to reproduce its
PIE~$\to$~JAAD-sized improvement here, but not confidently that it
makes this direction significantly worse. We report the point estimate
and its uncertainty together rather than rounding either away: the
fix's benefit is not confirmed to be architecture-independent, and a
practitioner should not assume it transfers to a new fusion design
without re-testing both directions at adequate $n$.

We probe both directions' cross-attention checkpoints
(Table~\ref{tab:probe-crossattn}), which reveals a more specific pattern
than "GRL + contrastive always increases separability." For
PIE~$\to$~JAAD, the direction that improves, separability rises
substantially over the unmodified baseline, both linearly
($0.844\pm0.006 \to 0.967\pm0.009$) and non-linearly
($0.846\pm0.000 \to 0.952\pm0.020$) -- matching every stacked-GRU
condition in Table~\ref{tab:probe}. For JAAD~$\to$~PIE, the direction
where GRL + contrastive fails to reproduce its improvement (its point
estimate falls, though not by more than the baseline's own noise;
Section~\ref{sec:results-crossattn}), separability barely moves at all
($0.845\pm0.004 \to 0.851\pm0.003$ linear; $0.828\pm0.033 \to
0.844\pm0.004$ non-linear) -- an order of magnitude smaller shift than
the PIE~$\to$~JAAD case. Separability and task performance are
therefore not simply decoupled in the sense of moving in opposite
directions; rather, the training procedure's measurable effect on the
representation tracks whether it helps the task at all. Where GRL +
contrastive meaningfully changes cross-dataset AUC (for better, on
stacked-GRU in both directions and cross-attention PIE~$\to$~JAAD), it
also meaningfully changes probe separability, always upward, never
downward; where it does not help (cross-attention JAAD~$\to$~PIE), it
leaves the representation close to where the unmodified baseline left
it, and the AUC change in that case is presumably driven by something
else in training (e.g.\ optimization dynamics specific to this
architecture and direction) rather than by any representational shift
this probe can detect. This refines, rather than weakens, our reading
in Section~\ref{sec:results-probe}: literal domain-invariance is not
what GRL + contrastive achieves in any condition we test, and the
magnitude of its (upward) effect on separability appears coupled to
whether it changes task behavior at all, not decoupled from it as a
simpler "always more separable" story would suggest.

\begin{table}[t]
\centering
\caption{Cross-attention probe accuracy, both transfer directions.
Baseline uses the checkpoint trained on the row's source dataset
(PIE-trained for the PIE$\to$JAAD row, JAAD-trained for JAAD$\to$PIE).}
\label{tab:probe-crossattn}
\begin{tabular}{lcc}
\toprule
Condition & Linear acc. & MLP acc. \\
\midrule
Baseline (PIE-trained) & $0.844\pm0.006$ & $0.846\pm0.000$ \\
GRL + contrastive, PIE$\to$JAAD & $0.967\pm0.009$ & $0.952\pm0.020$ \\
\midrule
Baseline (JAAD-trained) & $0.845\pm0.004$ & $0.828\pm0.033$ \\
GRL + contrastive, JAAD$\to$PIE & $0.851\pm0.003$ & $0.844\pm0.004$ \\
\bottomrule
\end{tabular}
\end{table}

\begin{table}[t]
\centering
\caption{Cross-attention architecture (box-anchored fusion), corrected
labels, $n=6$ seeds. Compare against Table~\ref{tab:fixes}'s
stacked-GRU numbers for the same comparison on SF-GRU.}
\label{tab:crossattn}
\begin{tabular}{lcc}
\toprule
Method & PIE$\to$JAAD & JAAD$\to$PIE \\
\midrule
Baseline (no adaptation) & $0.513\pm0.009$ & $0.512\pm0.188$ \\
GRL + contrastive & $0.559\pm0.027$ & $0.446\pm0.063$ \\
\bottomrule
\end{tabular}
\end{table}

\subsection{A Third Architecture: The Recipe Can Actively Hurt}
\label{sec:results-uncertainty}

We test a third, again structurally distinct architecture:
uncertainty-weighted fusion, in which each modality is independently
GRU-encoded and fused by a softmax over per-modality inverse-variance
weights (a learned confidence-weighted average) rather than sequential
concatenation (SF-GRU) or query-anchored attention (cross-attention).
This architecture is not an arbitrary third choice: it is the
strongest \emph{unadapted} baseline we find in an 8-way cross-dataset
architecture sweep (PIE~$\to$~JAAD AUC $0.585\pm0.009$, $n=3$,
exceeding SF-GRU's $0.498\pm0.009$ and cross-attention's
$0.513\pm0.009$ by a wide margin), making it the best-motivated
candidate for testing whether domain-adversarial training adds value
on top of an architecture that already generalizes comparatively well
without any adaptation. We apply the same GRL + contrastive recipe,
default hyperparameters, box-embedding contrastive tap on this
architecture's own independent box encoder output, $n=6$ seeds per
direction.

Table~\ref{tab:uncertainty} shows the recipe does not help here, and on
PIE~$\to$~JAAD actively \emph{hurts}: AUC falls from $0.585\pm0.009$ to
$0.499\pm0.072$, a drop that clears the two conditions' pooled standard
deviation and is the only case in this paper where GRL + contrastive
significantly \emph{under}performs an architecture's own unadapted
baseline. JAAD~$\to$~PIE also falls ($0.377\pm0.046 \to
0.351\pm0.060$), though this difference does not clear pooled standard
deviation. Combined with the cross-attention result
(Section~\ref{sec:results-crossattn}), only one of three architectures
we test (SF-GRU) shows the recipe helping in \emph{both} transfer
directions; cross-attention shows it helping in one and not
distinguishable from baseline in the other; uncertainty-weighted fusion
shows it significantly hurting in one and not distinguishable from
baseline (trending down) in the other. We revise our framing
accordingly: the earlier characterization of this recipe as producing
a ``partial, principled fix'' understates how architecture-contingent
it is. A more accurate summary is that GRL + contrastive is a
reproducible fix for SF-GRU's specific stacked-fusion topology, whose
mechanism (per Sections~\ref{sec:results-probe}
and~\ref{sec:results-crossattn}) does not appear to be literal
representational alignment, and which does not transfer as a
recommended recipe to the two alternative fusion mechanisms we
test -- one of which it actively harms on its strongest baseline.

\begin{table}[t]
\centering
\caption{Uncertainty-weighted fusion architecture, corrected labels.
Baseline is $n=3$ (from the 8-way architecture sweep); GRL + contrastive
is $n=6$. $\downarrow^{*}$ = drop clears pooled standard deviation;
$\downarrow$ = point estimate falls but not significantly at this $n$.}
\label{tab:uncertainty}
\begin{tabular}{lccc}
\toprule
Method & PIE$\to$JAAD & JAAD$\to$PIE & Verdict \\
\midrule
Baseline (no adaptation) & $0.585\pm0.009$ & $0.377\pm0.046$ & -- \\
GRL + contrastive & $0.499\pm0.072$ & $0.351\pm0.060$ & $\downarrow^{*}$/$\downarrow$ \\
\bottomrule
\end{tabular}
\end{table}

\subsection{A Pretrained Trajectory Representation Does Not Help}
\label{sec:results-pretrained-traj}

Gesnouin~et~al.~\cite{gesnouin2022} report that generic,
ImageNet/Sports1M-pretrained visual backbones generalize better under
PIE/JAAD domain shift than crossing-intention-specific visual
architectures -- motivating us to test the analogous question for the
\emph{motion} modality: does a box-trajectory representation pretrained
on a broader distribution than either dataset's own labeled training
pool generalize better cross-dataset than SF-GRU's from-scratch box
GRU? Lacking access to a large external trajectory corpus, we pretrain
a small Transformer encoder with a masked-frame reconstruction
objective (BERT-style: random frames of a box-delta sequence are
masked and predicted from context) on the pooled, \emph{unlabeled}
box-delta trajectories from both PIE's and JAAD's train and validation
splits combined -- a broader motion distribution than either dataset's
own labeled training pool alone, though a materially weaker substitute
for a genuine large-scale external corpus (e.g.\ Waymo Open Motion or
Argoverse~2), which we did not have local access to. We then freeze
this encoder and splice its per-timestep output into SF-GRU's box GRU
slot in place of the raw 4-dimensional box delta, training only the
downstream fusion (the box GRU, now consuming a 64-dimensional
embedding, and everything after it) under the same GRL + contrastive
recipe and default hyperparameters as Section~\ref{sec:results-dann}.

Table~\ref{tab:pretrained-traj} shows this does not help, and in fact
underperforms both the from-scratch GRL + contrastive result and, in
both directions, the plain unadapted baseline
(Table~\ref{tab:baseline}: $0.498$ and $0.402$). We do not read this as
evidence against a pretrained-backbone strategy in general, and we are
explicit about a confound in this experiment beyond corpus size alone:
our pretraining pool is drawn from PIE's and JAAD's \emph{own} train
and validation splits, the same two datasets used for the downstream
cross-dataset evaluation, whereas the visual-backbone analogy this
experiment is motivated by (Gesnouin~et~al.) concerns a backbone
pretrained on data with no relationship to either evaluation
domain. A self-supervised objective trained on a mixture of the source
and target domains has no pressure to discard whatever
dataset-identifying structure (frame rate, box-annotation convention,
scale) distinguishes the two -- indeed, our own probe methodology
(Section~\ref{sec:method-probe}) shows PIE and JAAD are separable from
nearly any representation of their inputs, including a single scalar
track length. It is therefore not clear whether this experiment tested
"does motion pretraining help cross-dataset transfer" or something
closer to "does an encoder pretrained to reconstruct a mixture of both
domains' own trajectories help," which is a different, and less
informative, question. Corpus size (a few hundred trajectories, vs.\
ImageNet's millions of images) compounds this: unlike a frozen VGG16
backbone pretrained on a large, source/target-independent corpus, a box
encoder pretrained on a small pool drawn from the very datasets under
study does not obviously carry information beyond what the
from-scratch, task-supervised box GRU already has access to, and the
added representation bottleneck (a frozen, non-task-tuned embedding
standing between the raw trajectory and the classification objective)
may cost more than the pretraining gives back at this scale. Isolating
which of these two factors -- domain composition or corpus size --
drives the negative result (for instance, by pretraining on PIE alone
and evaluating JAAD~$\to$~PIE, or vice versa, so the pretraining
corpus excludes the source domain entirely) is a direction for future
work we did not pursue here.

\begin{table}[t]
\centering
\caption{Frozen pretrained trajectory encoder spliced into the box
modality, GRL + contrastive recipe, default hyperparameters, corrected
labels, $n=6$ seeds. Compare against Table~\ref{tab:fixes}'s
from-scratch-box GRL + contrastive row (default hyperparameters:
$0.618\pm0.019$ / $0.583\pm0.073$).}
\label{tab:pretrained-traj}
\begin{tabular}{lcc}
\toprule
Method & PIE$\to$JAAD & JAAD$\to$PIE \\
\midrule
GRL + contrastive, pretrained traj.\ box & $0.489\pm0.019$ & $0.440\pm0.089$ \\
\bottomrule
\end{tabular}
\end{table}

\subsection{A Fine-Tuned Vision-Language Model Does Not Outperform SF-GRU}
\label{sec:results-vlm}

As a final, qualitatively different test of whether pretrained
world-knowledge helps this task, we LoRA-fine-tune Qwen2-VL-2B-Instruct
on a VQA-style (image-with-boxed-pedestrian, YES/NO) formulation of
crossing-intent, trained separately on each dataset and evaluated with
the same train-on-source/test-on-target protocol as every other result
in this paper. Unlike the from-scratch trajectory encoder of
Section~\ref{sec:results-pretrained-traj}, this backbone is pretrained
on a large, source/target-independent vision-language corpus, making it
a closer analogue to Gesnouin~et~al.'s generic-visual-backbone finding.
We report accuracy, precision, recall, and F1 rather than AUC: the
model is queried via greedy-decoded YES/NO generation, which yields a
hard prediction, not a calibrated probability, so these numbers are
\emph{not} directly comparable to the AUC figures used throughout the
rest of this paper.

In-domain accuracy is unremarkable relative to SF-GRU's in-domain AUC
(Table~\ref{tab:baseline}): $0.808$ (PIE, $n=224$) and, on JAAD, a
range of $0.647$--$0.824$ depending on which training checkpoint is
selected ($n=17$, see checkpoint-selection robustness below).
Cross-dataset, PIE~$\to$~JAAD accuracy is $0.581$ ($n=117$) with a
severe recall collapse (recall $0.154$: the model predicts
``not crossing'' for the large majority of samples regardless of
ground truth). For JAAD~$\to$~PIE, JAAD's validation split is only 17
sequences under \texttt{sample\_type=`beh'} -- the same structural
constraint noted throughout this paper for JAAD-as-source training
(Section~\ref{sec:results-discussion}) -- so selecting a single
``best'' fine-tuning checkpoint by validation accuracy on 17 samples is
itself high-variance (one flipped prediction moves accuracy by
$\approx\!6$ percentage points). Rather than report one
potentially-lucky checkpoint's cross-dataset number, we evaluate the
three best checkpoints by validation accuracy (training val.\ accuracy
$0.824$, $0.765$, $0.647$) independently on the full PIE test split
($n=637$ each). Table~\ref{tab:vlm-topk} shows this matters: the
checkpoint with the \emph{highest} training-time validation accuracy
is not the one with the best cross-dataset accuracy or F1, and the
spread across checkpoints (accuracy $0.281$--$0.474$) is large relative
to the differences we treat as meaningful elsewhere in this paper. We
report the mean and range across checkpoints rather than a single
point estimate for this reason.

Taken together, neither direction shows the fine-tuned VLM clearly
outperforming SF-GRU's from-scratch cross-dataset results (accuracy is
not directly comparable to our AUC numbers, but the severe recall
imbalance in several cells indicates the model is often not
discriminating well at all, independent of the choice of metric). We
read this as a second, qualitatively different negative result for the
``pretrained backbone rescues cross-dataset transfer'' hypothesis,
alongside Section~\ref{sec:results-pretrained-traj}'s trajectory-encoder
result -- though we note important differences in setup (a
single-image, VQA-style formulation with no explicit trajectory input,
vs.\ SF-GRU's multi-frame, multi-modality fusion) that mean this is not
a strictly controlled comparison, only a suggestive one.

\begin{table}[t]
\centering
\caption{VLM LoRA cross-dataset results, corrected labels. PIE$\to$JAAD
uses a single checkpoint (PIE's validation split, $n=224$, does not
have the small-sample selection problem JAAD's does). JAAD$\to$PIE
reports the three best JAAD-training checkpoints by validation accuracy
($n=17$ each), evaluated independently on PIE's full test split
($n=637$ each). Metrics are accuracy/F1/recall from greedy-decoded
YES/NO classification, not AUC.}
\label{tab:vlm-topk}
\begin{tabular}{lccc}
\toprule
Method & Acc. & F1 & Rec. \\
\midrule
PIE$\to$JAAD (single ckpt.) & $0.581$ & $0.246$ & $0.154$ \\
\midrule
JAAD$\to$PIE, rank 0 (best train val.) & $0.474$ & $0.394$ & $0.609$ \\
JAAD$\to$PIE, rank 1 & $0.422$ & $0.470$ & $0.911$ \\
JAAD$\to$PIE, rank 2 & $0.281$ & $0.439$ & $1.000$ \\
JAAD$\to$PIE, mean$\pm$std & $0.393\pm0.082$ & $0.434\pm0.031$ & $0.840\pm0.167$ \\
\bottomrule
\end{tabular}
\end{table}

\subsection{Split Sensitivity Is Not Explained by Video-Level Leakage}
\label{sec:results-leakage}

SF-GRU-family cross-dataset results are also sensitive to which
train/val/test partition is used within a dataset. Under JAAD's
official default split, our GRL + contrastive model reaches
PIE~$\to$~JAAD AUC $0.621\pm0.017$; under a random per-pedestrian
partition (\texttt{data\_split\_type=`random'}, ratios
$[0.5,0.4,0.1]$, freshly redrawn per run), the same method scores
$0.511\pm0.036$ -- a substantial drop. Because a per-pedestrian random
split does not guarantee that all pedestrians from a given video stay
in the same split, we hypothesized this drop was explained by
train/test video overlap: the default random split places $52\%$ of
JAAD pedestrians and (via a pedestrian-count coincidence) effectively
$100\%$ of PIE pedestrians in videos that appear on both sides of the
train/test boundary. We built a video-grouped variant of the random
split (Section~\ref{sec:method-probe}'s counterpart for splitting,
grouping all pedestrians in a video before partitioning, guaranteeing
zero video overlap) and re-ran the same method under it. Contrary to
our hypothesis, video-grouped-split AUC is \emph{lower}, not higher,
than the leaky pedestrian-level random split: $0.445\pm0.054$ ($n=3$)
vs.\ $0.511\pm0.036$. Video-level leakage is therefore not the
explanation for the default-split-vs-random-split gap; the true cause
remains open. We report this as a negative result rather than
discarding it, since it rules out what seemed the most likely
explanation and narrows the space of remaining hypotheses (e.g.\ the
default splits may differ systematically in scene composition or
pedestrian-behavior distribution in a way a random partition does not
control for).

\subsection{Summary of Everything Attempted}
\label{sec:results-summary}

Table~\ref{tab:summary} collects every fix attempt in this paper in one
place, since the individual results are otherwise spread across eight
subsections written around different hypotheses, with a verdict
computed against pooled standard deviation rather than by eye. Of the
thirteen non-baseline interventions, only the domain-adversarial family
on SF-GRU (GRL alone, GRL + contrastive full-representation, GRL +
contrastive context-targeted) improves PIE~$\to$~JAAD \emph{and}
JAAD~$\to$~PIE by more than pooled noise in both directions. The
geometry- and normalization-based fixes (rows 2--5) all significantly
\emph{hurt} PIE~$\to$~JAAD and are statistically indistinguishable
from baseline on JAAD~$\to$~PIE -- not the symmetric two-direction harm
a first read of the raw numbers might suggest, since JAAD~$\to$~PIE's
wide baseline variance (row 1; Section~\ref{sec:results-discussion})
absorbs differences that would be significant in the other direction.
The context-targeted variant (row 9) is statistically indistinguishable
from the full-representation recipe (row 8) in both directions
(Section~\ref{sec:results-context}) -- narrowing the domain-adversarial
pressure to the diagnosed leakage source loses nothing. Synthetic speed
and the pretrained trajectory representation (rows 11--12) both
significantly hurt \emph{both} directions relative to the tuned
headline recipe; mixup (row 10) is statistically indistinguishable from
it in both directions despite its lower point estimate. Applied to the
two alternative fusion architectures (rows 13--14), the recipe's record
is mixed-to-negative: on cross-attention it significantly helps
PIE~$\to$~JAAD and is statistically indistinguishable from its own
baseline (not significantly worse) on JAAD~$\to$~PIE; on
uncertainty-weighted fusion -- the strongest unadapted baseline of any
architecture we test -- it significantly \emph{hurts} PIE~$\to$~JAAD
and trends down (not significant) on JAAD~$\to$~PIE. Only SF-GRU shows
the recipe helping in both directions, whether applied broadly or
narrowly. Table~\ref{tab:summary} reports
AUC throughout and so excludes the VLM LoRA results
(Section~\ref{sec:results-vlm}, which use greedy-decoded accuracy/F1
rather than a probability score); those results are a further,
qualitatively different negative data point -- a large pretrained
vision-language backbone also does not clearly outperform SF-GRU's
from-scratch cross-dataset numbers -- that this table's AUC-only format
cannot represent alongside the rest. We read this table as the
paper's most concise evidence for its central claim: cross-dataset
collapse in this model family is not fixed by any single targeted
intervention we tried except domain-adversarial training on SF-GRU
specifically, and that fix is not just partial but architecture-
contingent, actively harmful on at least one alternative architecture's
strongest baseline.

\begin{table}[t]
\centering
\caption{Every intervention attempted in this paper, both transfer
directions, corrected labels. SF-GRU (stacked fusion) unless noted; row
4 (homography) is the metric ground-plane reprojection of
Section~\ref{sec:method-geometry}, row 5 drops the box modality
entirely. Verdict is $\uparrow$/$\downarrow$/$\approx$ relative to the
stated baseline, where $\approx$ means the difference is smaller than
the two conditions' pooled standard deviation (i.e.\ not
distinguishable at this $n$) -- rows 2--9 compare against the unadapted
baseline (row 1); rows 10--12 compare against GRL + contrastive tuned
(row 8); rows 13--14 compare against their own architecture's unadapted
baseline (Tables~\ref{tab:crossattn} and~\ref{tab:uncertainty}).}
\label{tab:summary}
\small
\begin{tabular}{lccc}
\toprule
Intervention & PIE$\to$JAAD & JAAD$\to$PIE & Verdict \\
\midrule
1. Baseline (no fix) & $0.498\pm0.009$ & $0.402\pm0.138$ & -- \\
2. Scale-norm & $0.437\pm0.025$ & $0.536\pm0.063$ & $\downarrow$/$\approx$ \\
3. Height-norm & $0.390\pm0.017$ & $0.554\pm0.104$ & $\downarrow$/$\approx$ \\
4. Homography & $0.406\pm0.011$ & $0.544\pm0.053$ & $\downarrow$/$\approx$ \\
5. No-box & $0.385\pm0.009$ & $0.432\pm0.085$ & $\downarrow$/$\approx$ \\
6. Contrastive only & $0.558\pm0.028$ & $0.450\pm0.107$ & $\uparrow$/$\approx$ \\
7. GRL only & $0.608\pm0.030$ & $0.587\pm0.069$ & $\uparrow$/$\uparrow$ \\
8. GRL+contrastive (tuned) & $\mathbf{0.621\pm0.017}$ & $\mathbf{0.613\pm0.071}$ & $\uparrow$/$\uparrow$ \\
9. GRL+contrastive (context-targeted) & $0.617\pm0.020$ & $0.637\pm0.055$ & $\uparrow$/$\uparrow$ \\
10.\ + mixup & $0.580\pm0.077$ & $0.537\pm0.079$ & $\approx$/$\approx$ \\
11.\ + synth.\ speed & $0.602\pm0.001$ & $0.430\pm0.072$ & $\downarrow$/$\downarrow$ \\
12.\ + pretrained traj.\ & $0.489\pm0.019$ & $0.440\pm0.089$ & $\downarrow$/$\downarrow$ \\
13.\ Cross-attn.\ & $0.559\pm0.027$ & $0.446\pm0.063$ & $\uparrow$/$\approx$ \\
14.\ Uncertainty fusion & $0.499\pm0.072$ & $0.351\pm0.060$ & $\downarrow$/$\approx$ \\
\bottomrule
\end{tabular}
\end{table}

\subsection{Discussion}
\label{sec:results-discussion}

Four results together shape our interpretation. First, geometry- and
normalization-based fixes are uniformly ineffective
(Section~\ref{sec:results-geometry}), ruling out the most
straightforward explanation (differing camera setups) for
cross-dataset collapse. Second, the modality-ablation result
(Section~\ref{sec:results-ablation}) and the probe results
(Section~\ref{sec:results-probe}) together point away from a
single-modality, single-cause account: visual context matters most to
zero out, yet the mechanism that helps most (GRL) does not work by
suppressing a detectable dataset signal -- it \emph{increases} that
signal's detectability while improving task performance. We read this
as evidence that GRL's practical benefit in this setting is closer to
a generic regularization effect on the fused representation (e.g.\
reducing reliance on dataset-specific spurious correlations the main
task loss alone would otherwise fit) than to the literal
representation-alignment mechanism it is designed around, echoing
similar decoupling results reported for DANN-family methods outside
this domain. Third, the cross-attention replication
(Section~\ref{sec:results-crossattn}) sharpens this reading further:
probing both transfer directions shows the separability increase is
not a fixed side effect of the training procedure, but tracks whether
GRL + contrastive changes task behavior at all -- large where it helps
(PIE~$\to$~JAAD on both architectures, JAAD~$\to$~PIE on stacked-GRU)
and negligible where it does not (JAAD~$\to$~PIE on cross-attention). A
representation-alignment account predicts a consistent effect on
separability regardless of whether AUC improves; what we observe -- the
separability shift and the AUC shift rising and falling together in
magnitude, even though the AUC shift is not always positive -- is more
consistent with GRL + contrastive perturbing the representation only
when it meaningfully engages with training at all, and that engagement
sometimes helping and sometimes hurting the task depending on
architecture-specific training dynamics. The uncertainty-weighted
fusion result (Section~\ref{sec:results-uncertainty}) is the sharpest
data point for this view: applied to the architecture with the
\emph{best} unadapted baseline of any we test, the same recipe
significantly \emph{degrades} it. If GRL + contrastive's benefit came
from a mechanism intrinsic to domain-adversarial training working
correctly, we would not expect it to reliably help one architecture
and reliably hurt another; a generic-regularizer account, where the
training procedure's effect on a given architecture's specific
optimization landscape can go either way, accommodates this more
naturally than an account in which the method's success is attributed
to achieving its nominal objective. Fourth, the persistent, large
JAAD~$\to$~PIE variance across
nearly every method in Tables~\ref{tab:fixes},~\ref{tab:crossattn},
and~\ref{tab:uncertainty}
deserves an explicit alternative explanation before we attribute it to
the benchmark's scene composition: JAAD's default-split training set
has 103 labeled sequences, against PIE's 796 -- an 8-fold difference --
so JAAD-as-source runs are trained on a much smaller, and therefore
individually noisier, sample than PIE-as-source runs, quite apart from
any property of \emph{which} scenes each dataset contains. A smaller
source training set producing a noisier decision boundary, and hence
higher-variance transfer to any target, is a mundane statistical
explanation that competes directly with a benchmark-composition
account, and our experiments do not distinguish between them: every
JAAD-as-source condition in this paper confounds ``JAAD'' with
``the smaller-training-set direction.'' The counter-hypothesis
split-sensitivity result (Section~\ref{sec:results-leakage}) is
consistent with either account. We flag this rather than resolve it: a
clean test would hold training-set size fixed (e.g.\ subsampling PIE's
source training set to JAAD's size) and check whether PIE~$\to$~JAAD
variance rises to match, which we did not run. Together, this and the
split-sensitivity result suggest that some of
this task's difficulty is a property of the benchmark itself --
scene composition, training-set size, or both -- rather than purely a
property of any model or training method evaluated on it. We view
robust cross-dataset pedestrian intent prediction as still substantially
unsolved, with our contribution best read as (a) a fix that is
reproducible and genuinely useful for one specific architecture
(SF-GRU's stacked fusion) but is not a general-purpose recipe --
tested on two further, structurally distinct architectures, it helps
partially on one and significantly hurts the other's best baseline in
at least one direction -- together with a mechanistic account
(Sections~\ref{sec:results-probe}--\ref{sec:results-uncertainty}) of
why its effect is architecture-contingent rather than a property of
whether it achieves genuine domain-invariance, and (b) a set of
ruled-out hypotheses -- geometry, video-level leakage, a missing speed
modality, and a locally-pretrained trajectory representation -- that we
hope narrows the search space for future work more than our single
positive result advances it.
